% --- Archon physics formalization source begin ---
% archon:physics
% archon:covers PhyXMiniProblems/problem_phyx_mini_0855.lean
% archon:source-report reports/phyx_mini/problem_phyx_mini_0855.source.json

\chapter{Physics problem phyx\_mini\_0855}
\label{ch:PhyXMiniProblems_problem_phyx_mini_0855}

\paragraph{Problem source.}
\#\# Question

What is the electric potential energy of the group of charges?

\#\# Answer choices

- A: A: 1.2 \textbackslash{}times 10\textasciicircum{}\{19\}J

- B: B: -1.2 \textbackslash{}times 10\textasciicircum{}\{19\}J

- C: C: 2.2 \textbackslash{}times 10\textasciicircum{}\{19\}J

- D: D: 0J

\#\# Figure caption (auxiliary; use the image as primary evidence)

The image shows a diagram with three positive charges, each represented as a red circle with a plus sign inside. The charges are positioned at the vertices of a right triangle. Here's a detailed description:

1. **Top Left Charge:**
   - Labeled as "2.0 nC."

2. **Top Right Charge:**
   - Labeled as "3.0 nC."
   - Positioned 3.0 cm horizontally to the right of the top left charge.

3. **Bottom Left Charge:**
   - Labeled as "3.0 nC."
   - Positioned 4.0 cm vertically below the top left charge.

4. **Distances:**
   - The horizontal distance between the top left and top right charges is noted as "3.0 cm."
   - The vertical distance between the top left and bottom left charges is noted as "4.0 cm."

The arrangement forms a right triangle with the top left charge at the right angle. Dotted lines connect the charges, indicating the distances between them.

\#\# Dataset metadata

Electrostatics; Physical Model Grounding Reasoning, Multi-Formula Reasoning

\paragraph{Recorded answer/context.}
D: D: 0J

\paragraph{Figure/image path.}
/root/proposal\_for\_physic/science-mango/phyx\_mini\_run/phyx\_data/test\_image/855.png

\paragraph{Formalization target.}
create a compiling Lean file with sorry bodies at `PhyXMiniProblems/problem\_phyx\_mini\_0855.lean`. The Lean declarations must preserve the physical quantities, dimensions or dimensional roles, figure labels, governing-law hypotheses, and final relation expressed by this problem.
Use Mathlib/Physlib names found through LeanExplore where available. If a physics API is missing, introduce faithful local abstractions rather than scalar placeholder aliases.

\begin{theorem}[Physics formalization target]
\label{thm:physics:phyx_mini_0855:target}
\lean{PhyXMiniProblems.ProblemPhyXMini0855.problem_phyx_mini_0855}
\uses{def:physics:phyx-mini-0855:phyxminiproblems-problemphyxmini0855-energyinjoules, def:physics:phyx-mini-0855:phyxminiproblems-problemphyxmini0855-threepointchargesetup, def:physics:phyx-mini-0855:phyxminiproblems-problemphyxmini0855-matchesproblemandprimaryfigure, def:physics:phyx-mini-0855:phyxminiproblems-problemphyxmini0855-hasphysicalthreechargeparameters, def:physics:phyx-mini-0855:phyxminiproblems-problemphyxmini0855-satisfiesdepictedseparationgeometry, def:physics:phyx-mini-0855:phyxminiproblems-problemphyxmini0855-usestextbookvacuumcoulombconstant, def:physics:phyx-mini-0855:phyxminiproblems-problemphyxmini0855-satisfiesthreepointchargeelectrostaticenergylaw, def:physics:phyx-mini-0855:phyxminiproblems-problemphyxmini0855-displayedenergyinjoules}
The assigned autoformalize agent should translate this physics problem into Lean declarations in the covered file, with theorem and lemma proof bodies written as `by sorry`.
\end{theorem}
\begin{proof}
This is an autoformalization task, not a proof task. Produce statements that can later be proved by the physics prover without weakening the physical model.
\end{proof}

% --- Archon declaration topology begin ---
\section{Declaration topology}
\begin{definition}[Charge Quantity]
  \label{def:physics:phyx-mini-0855:phyxminiproblems-problemphyxmini0855-chargequantity}
  \lean{PhyXMiniProblems.ProblemPhyXMini0855.ChargeQuantity}
  A signed, unit-independent electric charge
\end{definition}
\begin{proof}
  This is the declared model definition of Charge Quantity.
\end{proof}
\begin{definition}[Length Quantity]
  \label{def:physics:phyx-mini-0855:phyxminiproblems-problemphyxmini0855-lengthquantity}
  \lean{PhyXMiniProblems.ProblemPhyXMini0855.LengthQuantity}
  A nonnegative, unit-independent physical separation
\end{definition}
\begin{proof}
  This is the declared model definition of Length Quantity.
\end{proof}
\begin{definition}[coulomb Constant Dimension]
  \label{def:physics:phyx-mini-0855:phyxminiproblems-problemphyxmini0855-coulombconstantdimension}
  \lean{PhyXMiniProblems.ProblemPhyXMini0855.coulombConstantDimension}
  The physical dimension M L\textasciicircum{}3 T\textasciicircum{}-2 C\textasciicircum{}-2 of Coulomb's constant
\end{definition}
\begin{proof}
  This is the declared model definition of coulomb Constant Dimension.
\end{proof}
\begin{definition}[Coulomb Constant Quantity]
  \label{def:physics:phyx-mini-0855:phyxminiproblems-problemphyxmini0855-coulombconstantquantity}
  \lean{PhyXMiniProblems.ProblemPhyXMini0855.CoulombConstantQuantity}
  \uses{def:physics:phyx-mini-0855:phyxminiproblems-problemphyxmini0855-coulombconstantdimension}
  A nonnegative, unit-independent Coulomb constant
\end{definition}
\begin{proof}
  This is the declared model definition of Coulomb Constant Quantity.
\end{proof}
\begin{definition}[Figure Plane]
  \label{def:physics:phyx-mini-0855:phyxminiproblems-problemphyxmini0855-figureplane}
  \lean{PhyXMiniProblems.ProblemPhyXMini0855.FigurePlane}
  The affine plane used for the metre-coordinate readouts in the figure
\end{definition}
\begin{proof}
  This is the declared model definition of Figure Plane.
\end{proof}
\begin{definition}[charge In Coulombs]
  \label{def:physics:phyx-mini-0855:phyxminiproblems-problemphyxmini0855-chargeincoulombs}
  \lean{PhyXMiniProblems.ProblemPhyXMini0855.chargeInCoulombs}
  \uses{def:physics:phyx-mini-0855:phyxminiproblems-problemphyxmini0855-chargequantity}
  Read a physical charge in coherent SI coulombs
\end{definition}
\begin{proof}
  This is the declared model definition of charge In Coulombs.
\end{proof}
\begin{definition}[charge In Nanocoulombs]
  \label{def:physics:phyx-mini-0855:phyxminiproblems-problemphyxmini0855-chargeinnanocoulombs}
  \lean{PhyXMiniProblems.ProblemPhyXMini0855.chargeInNanocoulombs}
  \uses{def:physics:phyx-mini-0855:phyxminiproblems-problemphyxmini0855-chargequantity, def:physics:phyx-mini-0855:phyxminiproblems-problemphyxmini0855-chargeincoulombs}
  Read a physical charge in nanocoulombs
\end{definition}
\begin{proof}
  This is the declared model definition of charge In Nanocoulombs.
\end{proof}
\begin{definition}[separation In Meters]
  \label{def:physics:phyx-mini-0855:phyxminiproblems-problemphyxmini0855-separationinmeters}
  \lean{PhyXMiniProblems.ProblemPhyXMini0855.separationInMeters}
  \uses{def:physics:phyx-mini-0855:phyxminiproblems-problemphyxmini0855-lengthquantity}
  Read a physical separation in coherent SI metres
\end{definition}
\begin{proof}
  This is the declared model definition of separation In Meters.
\end{proof}
\begin{definition}[separation In Centimeters]
  \label{def:physics:phyx-mini-0855:phyxminiproblems-problemphyxmini0855-separationincentimeters}
  \lean{PhyXMiniProblems.ProblemPhyXMini0855.separationInCentimeters}
  \uses{def:physics:phyx-mini-0855:phyxminiproblems-problemphyxmini0855-lengthquantity, def:physics:phyx-mini-0855:phyxminiproblems-problemphyxmini0855-separationinmeters}
  Read a physical separation in centimetres
\end{definition}
\begin{proof}
  This is the declared model definition of separation In Centimeters.
\end{proof}
\begin{definition}[coulomb Constant In Newton Square Meters Per Square Coulomb]
  \label{def:physics:phyx-mini-0855:phyxminiproblems-problemphyxmini0855-coulombconstantinnewtonsquaremeterspersquarecoulomb}
  \lean{PhyXMiniProblems.ProblemPhyXMini0855.coulombConstantInNewtonSquareMetersPerSquareCoulomb}
  \uses{def:physics:phyx-mini-0855:phyxminiproblems-problemphyxmini0855-coulombconstantquantity}
  Read Coulomb's constant in N m\textasciicircum{}2 / C\textasciicircum{}2 = J m / C\textasciicircum{}2
\end{definition}
\begin{proof}
  This is the declared model definition of coulomb Constant In Newton Square Meters Per Square Coulomb.
\end{proof}
\begin{definition}[energy In Joules]
  \label{def:physics:phyx-mini-0855:phyxminiproblems-problemphyxmini0855-energyinjoules}
  \lean{PhyXMiniProblems.ProblemPhyXMini0855.energyInJoules}
  Read electrostatic potential energy in coherent SI joules
\end{definition}
\begin{proof}
  This is the declared model definition of energy In Joules.
\end{proof}
\begin{definition}[Charge Site]
  \label{def:physics:phyx-mini-0855:phyxminiproblems-problemphyxmini0855-chargesite}
  \lean{PhyXMiniProblems.ProblemPhyXMini0855.ChargeSite}
  The three charge locations identifiable in the supplied image
\end{definition}
\begin{proof}
  This is the declared model definition of Charge Site.
\end{proof}
\begin{definition}[Labelled Side]
  \label{def:physics:phyx-mini-0855:phyxminiproblems-problemphyxmini0855-labelledside}
  \lean{PhyXMiniProblems.ProblemPhyXMini0855.LabelledSide}
  The two dotted, dimension-labelled sides drawn in the image
\end{definition}
\begin{proof}
  This is the declared model definition of Labelled Side.
\end{proof}
\begin{definition}[Three Charge Figure]
  \label{def:physics:phyx-mini-0855:phyxminiproblems-problemphyxmini0855-threechargefigure}
  \lean{PhyXMiniProblems.ProblemPhyXMini0855.ThreeChargeFigure}
  \uses{def:physics:phyx-mini-0855:phyxminiproblems-problemphyxmini0855-chargesite, def:physics:phyx-mini-0855:phyxminiproblems-problemphyxmini0855-labelledside}
  Literal evidence carried by the primary raster. Its scalar labels are kept separate from the physical quantities in the setup
\end{definition}
\begin{proof}
  This is the declared model definition of Three Charge Figure.
\end{proof}
\begin{definition}[Ambient Medium]
  \label{def:physics:phyx-mini-0855:phyxminiproblems-problemphyxmini0855-ambientmedium}
  \lean{PhyXMiniProblems.ProblemPhyXMini0855.AmbientMedium}
  The material environment in which the electrostatic interaction occurs
\end{definition}
\begin{proof}
  This is the declared model definition of Ambient Medium.
\end{proof}
\begin{definition}[Potential Energy Reference]
  \label{def:physics:phyx-mini-0855:phyxminiproblems-problemphyxmini0855-potentialenergyreference}
  \lean{PhyXMiniProblems.ProblemPhyXMini0855.PotentialEnergyReference}
  Reference configuration used to assign electrostatic potential energy
\end{definition}
\begin{proof}
  This is the declared model definition of Potential Energy Reference.
\end{proof}
\begin{definition}[Three Point Charge Setup]
  \label{def:physics:phyx-mini-0855:phyxminiproblems-problemphyxmini0855-threepointchargesetup}
  \lean{PhyXMiniProblems.ProblemPhyXMini0855.ThreePointChargeSetup}
  \uses{def:physics:phyx-mini-0855:phyxminiproblems-problemphyxmini0855-chargequantity, def:physics:phyx-mini-0855:phyxminiproblems-problemphyxmini0855-lengthquantity, def:physics:phyx-mini-0855:phyxminiproblems-problemphyxmini0855-coulombconstantquantity, def:physics:phyx-mini-0855:phyxminiproblems-problemphyxmini0855-figureplane, def:physics:phyx-mini-0855:phyxminiproblems-problemphyxmini0855-chargesite, def:physics:phyx-mini-0855:phyxminiproblems-problemphyxmini0855-threechargefigure, def:physics:phyx-mini-0855:phyxminiproblems-problemphyxmini0855-ambientmedium, def:physics:phyx-mini-0855:phyxminiproblems-problemphyxmini0855-potentialenergyreference}
  Independent physical quantities for the three-charge configuration. The potential energy is an observable field and is not defined from the requested numeric answer
\end{definition}
\begin{proof}
  This is the declared model definition of Three Point Charge Setup.
\end{proof}
\begin{definition}[Matches Problem And Primary Figure]
  \label{def:physics:phyx-mini-0855:phyxminiproblems-problemphyxmini0855-matchesproblemandprimaryfigure}
  \lean{PhyXMiniProblems.ProblemPhyXMini0855.MatchesProblemAndPrimaryFigure}
  \uses{def:physics:phyx-mini-0855:phyxminiproblems-problemphyxmini0855-chargeinnanocoulombs, def:physics:phyx-mini-0855:phyxminiproblems-problemphyxmini0855-separationincentimeters, def:physics:phyx-mini-0855:phyxminiproblems-problemphyxmini0855-threepointchargesetup}
  Numerical and qualitative information read directly from image 855.png. The coordinate origin is chosen at the top-left charge, the positive horizontal axis points right, and the positive vertical axis points up. The unprinted diagonal distance and the electrostatic energy do not occur in this structure
\end{definition}
\begin{proof}
  This is the declared model definition of Matches Problem And Primary Figure.
\end{proof}
\begin{definition}[Has Physical Three Charge Parameters]
  \label{def:physics:phyx-mini-0855:phyxminiproblems-problemphyxmini0855-hasphysicalthreechargeparameters}
  \lean{PhyXMiniProblems.ProblemPhyXMini0855.HasPhysicalThreeChargeParameters}
  \uses{def:physics:phyx-mini-0855:phyxminiproblems-problemphyxmini0855-chargeincoulombs, def:physics:phyx-mini-0855:phyxminiproblems-problemphyxmini0855-separationinmeters, def:physics:phyx-mini-0855:phyxminiproblems-problemphyxmini0855-coulombconstantinnewtonsquaremeterspersquarecoulomb, def:physics:phyx-mini-0855:phyxminiproblems-problemphyxmini0855-threepointchargesetup}
  Positivity and nondegeneracy conditions for the physical configuration
\end{definition}
\begin{proof}
  This is the declared model definition of Has Physical Three Charge Parameters.
\end{proof}
\begin{definition}[Satisfies Depicted Separation Geometry]
  \label{def:physics:phyx-mini-0855:phyxminiproblems-problemphyxmini0855-satisfiesdepictedseparationgeometry}
  \lean{PhyXMiniProblems.ProblemPhyXMini0855.SatisfiesDepictedSeparationGeometry}
  \uses{def:physics:phyx-mini-0855:phyxminiproblems-problemphyxmini0855-separationinmeters, def:physics:phyx-mini-0855:phyxminiproblems-problemphyxmini0855-threepointchargesetup}
  Each dimensionful pair separation has the same metre readout as the metric distance between the corresponding plotted positions. This generic geometry law does not assume the unlabelled 5 cm diagonal
\end{definition}
\begin{proof}
  This is the declared model definition of Satisfies Depicted Separation Geometry.
\end{proof}
\begin{definition}[Uses Textbook Vacuum Coulomb Constant]
  \label{def:physics:phyx-mini-0855:phyxminiproblems-problemphyxmini0855-usestextbookvacuumcoulombconstant}
  \lean{PhyXMiniProblems.ProblemPhyXMini0855.UsesTextbookVacuumCoulombConstant}
  \uses{def:physics:phyx-mini-0855:phyxminiproblems-problemphyxmini0855-coulombconstantinnewtonsquaremeterspersquarecoulomb, def:physics:phyx-mini-0855:phyxminiproblems-problemphyxmini0855-threepointchargesetup}
  The dimensionful constant is calibrated against Physlib's definition 1 / (4 pi epsilon\_0), and that real readout uses the conventional 8.99 * 10\textasciicircum{}9 N m\textasciicircum{}2 / C\textasciicircum{}2 textbook approximation. This is universal calibration data, not a conclusion about the current charge configuration
\end{definition}
\begin{proof}
  This is the declared model definition of Uses Textbook Vacuum Coulomb Constant.
\end{proof}
\begin{definition}[Satisfies Three Point Charge Electrostatic Energy Law]
  \label{def:physics:phyx-mini-0855:phyxminiproblems-problemphyxmini0855-satisfiesthreepointchargeelectrostaticenergylaw}
  \lean{PhyXMiniProblems.ProblemPhyXMini0855.SatisfiesThreePointChargeElectrostaticEnergyLaw}
  \uses{def:physics:phyx-mini-0855:phyxminiproblems-problemphyxmini0855-chargeincoulombs, def:physics:phyx-mini-0855:phyxminiproblems-problemphyxmini0855-separationinmeters, def:physics:phyx-mini-0855:phyxminiproblems-problemphyxmini0855-coulombconstantinnewtonsquaremeterspersquarecoulomb, def:physics:phyx-mini-0855:phyxminiproblems-problemphyxmini0855-energyinjoules, def:physics:phyx-mini-0855:phyxminiproblems-problemphyxmini0855-threepointchargesetup}
  For point charges in vacuum with zero potential energy at infinite separation, the total energy is the sum over the three unordered pairs, k q\_i q\_j / r\_ij. The field is a governing law and contains no numeric answer or answer-choice label
\end{definition}
\begin{proof}
  This is the declared model definition of Satisfies Three Point Charge Electrostatic Energy Law.
\end{proof}
\begin{lemma}[diagonal Separation In Centimeters]
  \label{lem:physics:phyx-mini-0855:phyxminiproblems-problemphyxmini0855-diagonalseparationincentimeters}
  \lean{PhyXMiniProblems.ProblemPhyXMini0855.diagonalSeparationInCentimeters}
  \uses{def:physics:phyx-mini-0855:phyxminiproblems-problemphyxmini0855-separationincentimeters, def:physics:phyx-mini-0855:phyxminiproblems-problemphyxmini0855-threepointchargesetup, def:physics:phyx-mini-0855:phyxminiproblems-problemphyxmini0855-matchesproblemandprimaryfigure, def:physics:phyx-mini-0855:phyxminiproblems-problemphyxmini0855-satisfiesdepictedseparationgeometry}
  The unlabelled hypotenuse of the depicted 3-4-5 triangle is 5 cm
\end{lemma}
\begin{proof}
  The conclusion follows from the governing relations and figure readouts named in the statement of diagonal Separation In Centimeters.
\end{proof}
\begin{lemma}[pairwise Charge Distance Factor]
  \label{lem:physics:phyx-mini-0855:phyxminiproblems-problemphyxmini0855-pairwisechargedistancefactor}
  \lean{PhyXMiniProblems.ProblemPhyXMini0855.pairwiseChargeDistanceFactor}
  \uses{def:physics:phyx-mini-0855:phyxminiproblems-problemphyxmini0855-chargeincoulombs, def:physics:phyx-mini-0855:phyxminiproblems-problemphyxmini0855-separationinmeters, def:physics:phyx-mini-0855:phyxminiproblems-problemphyxmini0855-threepointchargesetup, def:physics:phyx-mini-0855:phyxminiproblems-problemphyxmini0855-matchesproblemandprimaryfigure, def:physics:phyx-mini-0855:phyxminiproblems-problemphyxmini0855-satisfiesdepictedseparationgeometry}
  After substituting the three charges and the 3, 4, 5 cm separations, the sum of q\_i q\_j / r\_ij is exactly 53 / 10\textasciicircum{}17 C\textasciicircum{}2/m
\end{lemma}
\begin{proof}
  The conclusion follows from the governing relations and figure readouts named in the statement of pairwise Charge Distance Factor.
\end{proof}
\begin{definition}[Answer Choice]
  \label{def:physics:phyx-mini-0855:phyxminiproblems-problemphyxmini0855-answerchoice}
  \lean{PhyXMiniProblems.ProblemPhyXMini0855.AnswerChoice}
  The answer labels supplied with the dataset item
\end{definition}
\begin{proof}
  This is the declared model definition of Answer Choice.
\end{proof}
\begin{definition}[displayed Energy In Joules]
  \label{def:physics:phyx-mini-0855:phyxminiproblems-problemphyxmini0855-displayedenergyinjoules}
  \lean{PhyXMiniProblems.ProblemPhyXMini0855.displayedEnergyInJoules}
  \uses{def:physics:phyx-mini-0855:phyxminiproblems-problemphyxmini0855-answerchoice}
  Literal energy in joules printed beside each answer label. These values are source metadata and do not constrain the independent physical energy field
\end{definition}
\begin{proof}
  This is the declared model definition of displayed Energy In Joules.
\end{proof}
\begin{definition}[recorded Dataset Answer]
  \label{def:physics:phyx-mini-0855:phyxminiproblems-problemphyxmini0855-recordeddatasetanswer}
  \lean{PhyXMiniProblems.ProblemPhyXMini0855.recordedDatasetAnswer}
  \uses{def:physics:phyx-mini-0855:phyxminiproblems-problemphyxmini0855-answerchoice}
  The answer label recorded in the supplied dataset metadata
\end{definition}
\begin{proof}
  This is the declared model definition of recorded Dataset Answer.
\end{proof}
% --- Archon declaration topology end ---
% --- Archon physics formalization source end ---
